\documentclass[a4paper, 12pt]{article}

% ========== PACOTES ==========
\usepackage[utf8]{inputenc}
\usepackage[T1]{fontenc}
\usepackage[brazilian]{babel}
\usepackage{amsmath, amssymb}
\usepackage{booktabs}
\usepackage{geometry}
\usepackage{graphicx}
\usepackage{float}
\usepackage{caption}
\usepackage{xcolor}
\usepackage{hyperref}
\usepackage{fancyhdr}
\usepackage{lastpage}
\usepackage{siunitx}
\usepackage{multirow}
\usepackage{longtable}

\geometry{margin=2cm}
\sisetup{
    group-separator = {.},
    output-decimal-marker = {,},
    round-mode = places,
}

% ========== CORES ==========
\definecolor{verde}{RGB}{0, 128, 0}
\definecolor{vermelho}{RGB}{200, 0, 0}
\definecolor{azulescuro}{RGB}{0, 51, 102}

% ========== CABEÇALHO ==========
\pagestyle{fancy}
\fancyhf{}
\fancyhead[L]{\textcolor{azulescuro}{\textbf{QuantNucleo} — Backtest Cego 6M}}
\fancyhead[R]{\textcolor{azulescuro}{LMF / UFU}}
\fancyfoot[C]{Página \thepage\ de \pageref{LastPage}}
\renewcommand{\headrulewidth}{1pt}
\renewcommand{\footrulewidth}{0.4pt}

\hypersetup{
    colorlinks=true,
    linkcolor=azulescuro,
    urlcolor=azulescuro,
}

\begin{document}

% ========== CAPA ==========
\begin{titlepage}
\centering
\vspace*{3cm}
{\Huge\bfseries\textcolor{azulescuro}{Relatório de Backtest Cego}\\[0.5cm]}
{\LARGE Teste Out-of-Sample — Últimos 6 Meses\\[1cm]}
{\Large Alocação Fictícia de R\$ 10.000.000,00\\[2cm]}
{\large
\textbf{Núcleo Quantitativo — LMF / UFU}\\
Universidade Federal de Uberlândia\\[1cm]
Período: 30/08/2025 a 01/03/2026\\
Gerado em: 01/03/2026 00:54
}
\vfill
{\footnotesize Documento gerado automaticamente pelo QuantNucleo Backtesting Engine v0.1.0}
\end{titlepage}

\tableofcontents
\newpage

% ========================================================================
\section{Introdução e Metodologia}
% ========================================================================

Este relatório apresenta os resultados do \textbf{backtest cego (out-of-sample)} realizado
com uma alocação fictícia de \textbf{R\$ 10.000.000,00} no período de \textbf{30/08/2025} a \textbf{01/03/2026} (123 dias úteis).

\subsection{Objetivos}
\begin{itemize}
    \item Avaliar o desempenho de múltiplas estratégias quantitativas em um \textbf{teste cego}
          (dados não vistos durante o desenvolvimento);
    \item Comparar rentabilidade absoluta e ajustada ao risco;
    \item Mensurar o resultado de uma carteira diversificada (Renda Variável + Renda Fixa);
    \item Trazer os valores finais ao \textbf{valor presente} utilizando a taxa Selic/CDI como desconto.
\end{itemize}

\subsection{Alocação da Carteira}
A alocação entre classes de ativos segue a distribuição:
\begin{itemize}
    \item \textbf{70\%} em Renda Variável (ações B3):
          PETR4, VALE3, ITUB4, BBDC4, ABEV3, WEGE3, BBAS3, SUZB3, B3SA3, MGLU3;
    \item \textbf{30\%} em Renda Fixa (Tesouro Direto):
          LTN 2026, NTN-B 2030, NTN-F 2029.
\end{itemize}

\subsection{Estratégias de Renda Variável Avaliadas}
\begin{enumerate}
    \item \textbf{Buy \& Hold Equal-Weight}: compra e manutenção com pesos iguais (benchmark passivo);
    \item \textbf{Momentum Crossover MM20/MM50}: sinal de compra quando a média móvel de 20 dias
          cruza acima da de 50 dias; zera posição no cruzamento inverso;
    \item \textbf{Risk-Weighted (Vol Inversa)}: pesos proporcionais ao inverso da volatilidade
          realizada (63 dias), rebalanceados diariamente.
\end{enumerate}

\subsection{Benchmark}
O benchmark de referência é o \textbf{CDI acumulado} no período, equivalente a uma
taxa de 13.75\% a.a. O CDI acumulado no período foi de:

$$\text{CDI}_{acum} = (1 + 0.00051137)^{123} - 1 = 6.4902\%$$

$$\text{Valor CDI} = R\$ 10,649,019.31$$

% ========================================================================
\newpage
\section{Resultados — Renda Variável}
% ========================================================================

\subsection{Comparativo de Performance}

\begin{table}[H]
\centering
\caption{Comparativo de performance das estratégias de Renda Variável}
\label{tab:comparativo_rv}
\resizebox{\textwidth}{!}{%
    \begin{tabular}{lccccccc}
    \toprule
    \textbf{Estratégia} & \textbf{Valor Final} & \textbf{Retorno} & \textbf{Ret. Anual.} &
    \textbf{Vol. Anual.} & \textbf{Sharpe} & \textbf{Sortino} & \textbf{Max DD} \\
    \midrule
        Buy \& Hold Equal-Weight & R\$ 9,537,295.61 & 36.25\% & 89.44\% & 16.59\% & 4.5622 & 7.9511 & -5.83\% \\
        Momentum MM20/MM50 & R\$ 8,780,680.71 & 25.44\% & 59.10\% & 16.82\% & 2.6953 & 4.3225 & -6.62\% \\
        Risk-Weighted (Vol Inversa) & R\$ 9,630,213.42 & 37.57\% & 92.24\% & 14.83\% & 5.2931 & 8.9109 & -4.58\% \\
    \bottomrule
    \end{tabular}%
}
\end{table}

\subsection{VaR, CVaR e Estatísticas de Trading}

\begin{table}[H]
\centering
\caption{Value at Risk, CVaR e estatísticas operacionais}
\label{tab:var_cvar}
    \begin{tabular}{lcccc}
    \toprule
    \textbf{Estratégia} & \textbf{VaR 95\% Diário} & \textbf{CVaR 95\%} & \textbf{Win Rate} & \textbf{Trades} \\
    \midrule
        Buy \& Hold Equal-Weight & -0.9499\% & -1.6726\% & 59.8\% & N/A \\
        Momentum MM20/MM50 & -1.2759\% & -1.9285\% & 54.5\% & 28 \\
        Risk-Weighted (Vol Inversa) & -0.8872\% & -1.5399\% & 59.4\% & N/A \\
    \bottomrule
    \end{tabular}
\end{table}

\subsection{Interpretação}
\begin{itemize}
    \item O \textbf{Sharpe Ratio} é calculado como $\frac{R_{anual} - CDI_{anual}}{\sigma_{anual}}$,
          portanto valores positivos indicam retorno excedente ao CDI;
    \item O \textbf{Sortino Ratio} penaliza apenas a volatilidade negativa (downside deviation);
    \item O \textbf{Max Drawdown} é a maior queda pico-a-vale no período;
    \item O \textbf{VaR 95\%} indica a perda máxima esperada em 95\% dos dias.
\end{itemize}

% ========================================================================
\newpage
\section{Resultados — Renda Fixa}
% ========================================================================

\subsection{Composição da Carteira RF}
Capital alocado em renda fixa: \textbf{R\$ 3,000,000.00}.

\begin{table}[H]
\centering
\caption{Detalhamento dos títulos de renda fixa e marcação a mercado}
\label{tab:rf_detalhe}
\resizebox{\textwidth}{!}{%
    \begin{tabular}{lccccccc}
    \toprule
    \textbf{Título} & \textbf{Peso} & \textbf{Capital} & \textbf{PU Ini} & \textbf{PU Fim} &
    \textbf{Retorno} & \textbf{Ganho R\$} & \textbf{$\Delta$ Taxa (bps)} \\
    \midrule
        LTN 2026 (Prefixado) & 40\% & R\$ 1,200,000.00 & 893.26 & 946.24 & 5.93\% & R\$ 71,172.80 & -55 \\
        NTN-B 2030 (IPCA+) & 30\% & R\$ 900,000.00 & 743.74 & 773.74 & 4.03\% & R\$ 36,303.42 & -25 \\
        NTN-F 2029 (Prefixado c/ cupom) & 30\% & R\$ 900,000.00 & 706.09 & 755.65 & 7.02\% & R\$ 63,164.75 & -50 \\
    \midrule
    \textbf{Total RF} & 100\% &
    \textbf{R\$ 3,000,000.00} & — & — &
    \textbf{5.69\%} &
    \textbf{R\$ 170,640.97} & — \\
    \bottomrule
    \end{tabular}%
}
\end{table}

\subsection{Interpretação MtM}
A marcação a mercado mostra que a \textbf{queda nas taxas de juros} no período gerou
ganho de capital nos títulos prefixados e indexados ao IPCA. O Delta Taxa (em basis points)
indica a variação entre a taxa de compra e a taxa de mercado na data de referência.

A Duration Modificada amplifica o impacto: para cada $-100$ bps de queda na taxa,
o PU sobe aproximadamente $D_{mod} \times 1\%$.

% ========================================================================
\newpage
\section{Carteira Consolidada e Valor Presente}
% ========================================================================

\subsection{Consolidação RV + RF}

\begin{table}[H]
\centering
\caption{Resultado consolidado da carteira}
\label{tab:consolidado}
    \begin{tabular}{lc}
    \toprule
    \textbf{Métrica} & \textbf{Valor} \\
    \midrule
    Capital Inicial Total & R\$ 10,000,000.00 \\
    Capital RV (70\%) & R\$ 7,000,000.00 \\
    Capital RF (30\%) & R\$ 3,000,000.00 \\
    \midrule
    Valor Final RV & R\$ 9,630,215.00 \\
    Valor Final RF & R\$ 3,170,640.97 \\
    \midrule
    \textbf{Valor Final Total} & \textbf{R\$ 12,800,855.97} \\
    \textbf{Retorno Total} & \textbf{28.01\%} \\
    \textbf{Retorno Anualizado} & \textbf{65.85\%} \\
    \midrule
    CDI Acumulado (benchmark) & 6.49\% \\
    Valor se 100\% CDI & R\$ 10,649,019.31 \\
    Alpha (vs CDI) & 21.52\% \\
    \bottomrule
    \end{tabular}
\end{table}

\subsection{Valor Presente (Desconto a CDI)}

Todos os valores finais são trazidos ao valor presente pela taxa CDI de
13.75\% a.a., descontando 123 dias úteis:

$$VP = \frac{VF}{(1 + CDI_{a.a.})^{du / 252}}$$

\begin{table}[H]
\centering
\caption{Valores finais trazidos ao presente (desconto a CDI)}
\label{tab:valor_presente}
    \begin{tabular}{lccc}
    \toprule
    \textbf{Estratégia} & \textbf{Valor Final (VF)} & \textbf{Valor Presente (VP)} & \textbf{VP/Capital - 1} \\
    \midrule
        Buy \& Hold Equal-Weight (RV) & R\$ 9,537,295.61 & R\$ 8,956,031.85 & -10.44\% \\
        Momentum MM20/MM50 (RV) & R\$ 8,780,680.71 & R\$ 8,245,529.90 & -17.54\% \\
        Risk-Weighted (Vol Inversa) (RV) & R\$ 9,630,213.42 & R\$ 9,043,286.65 & -9.57\% \\
        Renda Fixa & R\$ 3,170,640.97 & R\$ 2,977,401.84 & -0.75\% \\
        \midrule
        \textbf{Consolidada} & \textbf{R\$ 12,800,855.97} & \textbf{R\$ 12,020,689.98} & \textbf{20.21\%} \\
    \bottomrule
    \end{tabular}
\end{table}

\subsection{Interpretação do Valor Presente}
Ao trazer os valores ao presente, eliminamos o efeito do ``custo de oportunidade'' do CDI.
Um VP maior que o capital inicial indica que a estratégia gerou \textbf{alpha real} acima do
benchmark risk-free. Valores presentes negativos (VP $<$ Capital) indicam destruição de valor
em termos reais.

% ========================================================================
\newpage
\section{Conclusão}
% ========================================================================

O backtest cego dos últimos 6 meses, com alocação fictícia de R\$ 10 milhões,
permite as seguintes conclusões:

\begin{enumerate}
    \item \textbf{Diversificação RV + RF}: a carteira consolidada (70\% RV + 30\% RF)
          obteve retorno de 28.01\%, superando o CDI acumulado de 6.49\%;

    \item \textbf{Renda Fixa}: o componente RF contribuiu com retorno de 5.69\%,
          beneficiado pela queda nas taxas de juros
          e consequente ganho de marcação a mercado;

    \item \textbf{Risk-Weighted}: a estratégia de pesos por volatilidade inversa tende a apresentar
          menor drawdown, sendo adequada para investidores conservadores;

    \item \textbf{Momentum}: a estratégia de cruzamento de médias apresenta desempenho variável
          dependendo do regime de mercado (tendencial vs lateral);

    \item \textbf{Valor Presente}: ao descontar pelo CDI, o VP consolidado resultou em
          R\$ 12,020,689.98, indicando geração de alpha real.
\end{enumerate}

\vspace{1cm}
\noindent\rule{\textwidth}{0.4pt}
\begin{center}
\textit{Este relatório é meramente informativo e acadêmico. Não constitui recomendação de investimento.\\
Resultados passados não garantem retornos futuros.}
\end{center}

\end{document}
